\documentclass[11pt]{article}
\usepackage[margin=1in]{geometry}
\usepackage{amsmath,amssymb,amsthm,mathtools}
\usepackage{graphicx}
\usepackage{microtype}
\usepackage[hidelinks]{hyperref}

\newtheorem{theorem}{Theorem}
\newtheorem{lemma}{Lemma}
\newtheorem{corollary}{Corollary}
\newcommand{\R}{\mathbb R}
\newcommand{\e}{\mathrm e}

\title{A Nonconvex Gauge with a Global Thy Angle\\
\large Candidate full affirmative answers to Problems 1 and 2 of arXiv:0902.2731}
\author{Autonomous proof candidate for human review}
\date{11 August 2026}

\begin{document}
\maketitle

\begin{center}
\fbox{\parbox{0.91\textwidth}{\textbf{Status.}
Candidate full affirmative answers, likely valid.  An explicit
positive-definite, absolutely homogeneous gauge on $\R^2$, smooth away from
the origin, fails the triangle
inequality, while its normalized polarization product satisfies the
Cauchy--Schwarz--Bunjakowsky inequality.  Its Thy angle also satisfies axiom
(An 11).  The only large algebraic step is accompanied by an exact rational
Bernstein-basis certificate generated from the gauge itself.}}
\end{center}

\section{The questions and the answer}

For a positive-definite absolutely homogeneous gauge $q$, the source defines
the polarization product, for nonzero $x,y$, by
\begin{equation}\label{eq:polarization}
 \langle x\mid y\rangle_{\spadesuit}
 =\frac14q(x)q(y)\left[
 q\!\left(\frac{x}{q(x)}+\frac{y}{q(y)}\right)^2-
 q\!\left(\frac{x}{q(x)}-\frac{y}{q(y)}\right)^2
 \right].
\end{equation}
Problem 1 asks whether such a $q$ can fail the triangle inequality while
\(
 |\langle x\mid y\rangle_{\spadesuit}|\le q(x)q(y)
\)
for all $x,y$.  Problem 2 asks whether, in addition, for every independent
$x,y$ the map
\[
 t\longmapsto \angle_{\rm Thy}(x,y+tx)
 =\arccos\frac{\langle x\mid y+tx\rangle_{\spadesuit}}
 {q(x)q(y+tx)}
\]
is a decreasing homeomorphism from $\R$ onto $(0,\pi)$.

\begin{figure}[ht]
 \centering
 \includegraphics[width=0.98\textwidth]{figures/problem_crop.png}
 \caption{Problems 1 and 2 in the source, PDF page 19.}
\end{figure}

\begin{theorem}\label{thm:main}
Both problems have affirmative answers.  On $\R^2$, put $q(0)=0$ and, in
Euclidean polar coordinates,
\begin{equation}\label{eq:qpolar}
 q(r\cos\theta,r\sin\theta)
 =r\left(1+\frac1{10}\cos4\theta\right).
\end{equation}
Then $q$ is continuous, positive definite, and absolutely homogeneous, but
does not satisfy the triangle inequality.  Nevertheless its product
\eqref{eq:polarization} satisfies CSB, and the associated Thy angle satisfies
(An 11).
\end{theorem}

The rest of the packet proves the theorem.  The conceptual point is that the
small fourfold-symmetric inward dents are compatible with strict angular
monotonicity of the polarization, even though corners of the kind excluded in
the source are not.

\section{Elementary properties and failure of convexity}

Write
\[
 w(\theta)=1+\frac1{10}\cos4\theta,
 \qquad
 u_\theta=\frac{(\cos\theta,\sin\theta)}{w(\theta)}.
\]
Since $9/10\le w\le11/10$, formula \eqref{eq:qpolar} gives a continuous,
positive-definite, absolutely homogeneous gauge, and $q(u_\theta)=1$.  Away
from the origin it also has the Cartesian formula
\begin{equation}\label{eq:qcartesian}
 q(X,Y)^2=
 \frac{(11X^4+14X^2Y^2+11Y^4)^2}
 {100(X^2+Y^2)^3}.
\end{equation}

Let $h=\pi/8$.  Because $w(h)=w(-h)=1$, the two vectors
\[
 u_+=(\cos h,\sin h),\qquad u_-=(\cos h,-\sin h)
\]
have $q(u_+)=q(u_-)=1$.  Their midpoint is $(\cos h,0)$, but
\[
 q\!\left(\frac{u_++u_-}{2}\right)
 =\frac{11}{10}\cos\frac\pi8>1.
\]
Thus the unit ball is not convex.  More directly, the triangle inequality
and homogeneity would force the left side to be at most
$\tfrac12(q(u_+)+q(u_-))=1$, a contradiction.

\section{Reduction to strict angular monotonicity}

For a base direction $\alpha$ and an oriented separation $\delta$, define
\begin{equation}\label{eq:Cdef}
 C_\alpha(\delta)
 =\frac14\left(q(u_\alpha+u_{\alpha+\delta})^2
                -q(u_\alpha-u_{\alpha+\delta})^2\right).
\end{equation}
This is exactly the normalized product
$\langle u_\alpha\mid u_{\alpha+\delta}\rangle_\spadesuit$.

\begin{lemma}\label{lem:monotone}
For every $\alpha\in\R$, the function $C_\alpha$ is strictly decreasing on
$[0,\pi]$, with
\[
 C_\alpha(0)=1,\qquad C_\alpha(\pi)=-1.
\]
\end{lemma}

Before proving the lemma, note its two immediate consequences.  Any pair of
directions can, after exchanging the two vectors if necessary, be written
with a separation in $[0,\pi]$.  Lemma~\ref{lem:monotone} then gives
$|C_\alpha(\delta)|\le1$.  Absolute homogeneity of
\eqref{eq:polarization} therefore gives CSB for all vectors.  It remains only
to establish the strict derivative assertion.

\section{Exact positivity certificate}

The gauge is invariant under rotation through $\pi/2$ and reflection in the
first coordinate axis.  For $0<\delta\le\pi/2$, rotation therefore reduces
the derivative calculation to
\[
 0\le\alpha\le\frac\pi2,\qquad
 0<\delta\le\frac\pi2.
\]
Introduce rational half-angle coordinates
\[
 x=\tan\frac\alpha2\in[0,1],\qquad
 z=\tan\frac\delta2\in[0,1].
\]
For reference, a $q$-unit direction with half-angle parameter $t$ is the
rational vector
\begin{equation}\label{eq:unitrational}
 U(t)=\frac{10(1+t^2)^3}
 {11+12t^2+130t^4+12t^6+11t^8}\,(1-t^2,2t).
\end{equation}
Thus $u_\alpha=U(x)$ and, where the displayed fraction is finite,
\[
 u_{\alpha+\delta}=U\!\left(\frac{x+z}{1-xz}\right).
\]
Equations \eqref{eq:qcartesian} and \eqref{eq:unitrational} make
$C_\alpha(\delta)$ a rational function of $x,z$.

The accompanying verifier performs the following calculation from scratch
over $\mathbb Q[x,z]$.  It forms the two rational vectors $U(x)\pm
U((x+z)/(1-xz))$, evaluates \eqref{eq:qcartesian}, forms
\eqref{eq:Cdef}, differentiates, and cancels the polynomial gcd.  The result
has the form
\begin{equation}\label{eq:certificate}
 -\frac{1}{\sin\delta}\frac{\partial C_\alpha}{\partial\delta}
 =\frac{(1+z^2)^4P(x,z)}{H(x,z)},
\end{equation}
where $P,H$ are primitive integer polynomials.  (The primitive contents of
the reduced numerator and denominator are exactly equal.)  The identity first holds on
the dense locus on which the half-angle expressions have no pole and then,
after cancellation, by continuity throughout the relevant interior.

Here is the exact positivity certificate:
\begin{center}
\begin{tabular}{c|c|c|c|c}
 polynomial & bidegree & monomials & Bernstein coefficients & least coefficient\\ \hline
 $P$ & $(128,104)$ & $6773$ & $13545$ & $12531744508246953$\\
 $H$ & $(128,112)$ & $7289$ & $14577$ & $1285453186679270785/28$
\end{tabular}
\end{center}
Every tensor Bernstein coefficient on $[0,1]^2$ is strictly positive.  The
least coefficient of $P$ has tensor index $(0,0)$; that of $H$ has index
$(1,1)$.  Since Bernstein basis functions are nonnegative and sum to one,
this proves $P>0$ and $H>0$ on the whole square.  Equation
\eqref{eq:certificate} consequently gives
\begin{equation}\label{eq:firsthalf}
 \frac{\partial C_\alpha}{\partial\delta}<0
 \qquad(0<\delta\le\pi/2).
\end{equation}
The verifier uses Python's exact \texttt{Fraction} arithmetic for the entire
Bernstein conversion.  It stores no precomputed $P$ or $H$; those
polynomials are reconstructed from \eqref{eq:qpolar} on every run.

\begin{proof}[Completion of Lemma~\ref{lem:monotone}]
It remains to transfer \eqref{eq:firsthalf} to the second half-circle.  Put
$\varepsilon=\pi-\delta$ when $\pi/2<\delta<\pi$.  Antipodality, oddness of
the normalized polarization in either argument, and reflection invariance
give
\begin{align*}
 C_\alpha(\pi-\varepsilon)
 &=-C(u_\alpha,u_{\alpha-\varepsilon})\\
 &=-C(u_{-\alpha},u_{-\alpha+\varepsilon})
 =-C_{-\alpha}(\varepsilon).
\end{align*}
Since $d\varepsilon/d\delta=-1$, differentiation yields
\[
 \frac{d}{d\delta}C_\alpha(\delta)
 =C_{-\alpha}'(\varepsilon)<0
\]
by \eqref{eq:firsthalf}.  Finally, absolute homogeneity gives
$u_{\alpha+\pi}=-u_\alpha$, so direct substitution into
\eqref{eq:Cdef} gives the endpoint values $1$ and $-1$.
\end{proof}

\section{Axiom (An 11)}

Let $a,b\in\R^2$ be linearly independent.  Suppose first that $b$ lies in the
open counterclockwise half-plane from the ray $a$.  The counterclockwise
separation $\delta(t)$ from $a$ to $b+ta$ is a decreasing homeomorphism
\[
 \delta:\R\longrightarrow(0,\pi),
 \qquad
 \lim_{t\to-\infty}\delta(t)=\pi,\qquad
 \lim_{t\to+\infty}\delta(t)=0.
\]
(After a Euclidean choice of direction coordinates with the ray $a$ as the
positive horizontal axis, this is simply the argument of $(s+t,r)$ for some
$r>0$.)  Lemma~\ref{lem:monotone} says that the normalized
polarization cosine $C(\delta(t))$ is a strictly increasing homeomorphism
from $\R$ onto $(-1,1)$.  Since $\arccos$ is a decreasing homeomorphism from
$(-1,1)$ onto $(0,\pi)$,
\[
 t\longmapsto\angle_{\rm Thy}(a,b+ta)
\]
is a decreasing homeomorphism from $\R$ onto $(0,\pi)$.  If $b$ lies in the
clockwise half-plane, then
$C(u_\alpha,u_{\alpha-\delta})=C(u_{-\alpha},u_{-\alpha+\delta})$
by reflection invariance, and the identical argument applies with base
$-\alpha$.  This is exactly (An 11), and completes the proof of
Theorem~\ref{thm:main}.

\section{Verification and scope}

The exact verifier is \texttt{verify\_certificate.py} in this packet.  On
the recorded run it reconstructed the rational Thy cosine with numerator and
denominator bidegree $(96,84)$, canceled a derivative gcd of bidegree
$(64,57)$, and certified all $28122$ Bernstein coefficients described above.
The nonconvexity witness and the symmetry reduction are independent of that
calculation.

A bounded novelty search found the two questions in the 2009 source, the
author's later 2013 classification paper (which continues to discuss
nonconvex balanced weights and concave-corner obstructions), and later work
using Thy angles in normed or $n$-normed spaces.  It did not locate this
gauge, this certificate, or a later solution of the two questions.  Priority
is not asserted; the result should be independently reviewed before citation.

\begin{thebibliography}{9}
\bibitem{Thurey2009}
V.~W. Th\"urey,
\emph{Angles and Polar Coordinates In Real Normed Spaces},
arXiv:0902.2731 (2009).

\bibitem{Thurey2013}
V.~W. Th\"urey,
\emph{Angles and a Classification of Normed Spaces},
Annals of Functional Analysis 4 (2013), no.~1, 114--137;
arXiv:1207.0074.
\end{thebibliography}

\end{document}
